\section[Introduction]{Introduction}

\subsection[Course Organization]{Course Organization}

\begin{frame}{Course Organization}
	The course is organized around:
	\begin{itemize}
		\item Lectures: 20 hours (2h x 10 sessions)
		\item Tutorials: 16 hours (2h x 8 sessions)
		\item Practical sessions: 20 hours (2h x 10 sessions)
	\end{itemize}
	\vspace{0.5cm}
	The following time slots will be used:
	\begin{itemize}
		\item Tuesday 15:45 - 17:45: Tutorials / Practical sessions
		\item Tuesday 18:00 - 20:00: Exams
		\item Thursday 13:30 - 15:30: Lectures / Practical sessions
	\end{itemize}
\end{frame}

\begin{frame}{Assessments}
	\begin{tabular}{|l|l|l|l|l|}
		\hline
		\textbf{Item} & \textbf{\%} & \textbf{Duration} & \textbf{Format} & \textbf{Date} \\
		\hline
		CC1 & 33\% & 20 min & In-class, MCQ & 19/03 \\
		\hline
		TD \& TP & 33\% & - & Attendance, lab completion, project & - \\
		\hline
		CC3 & 34\% & 1h30 & Exam conditions & 19/05 \\
		\hline
		CC4 & - & 1h30 & Exam conditions & 11/06 \\
		\hline
	\end{tabular} \\
	\vspace{0.5cm}
	Note: The CC3 exam will take place during the Tuesday 18:00 - 20:00 slot.
\end{frame}

\begin{frame}{Course Outline}
	\begin{itemize}
		\item Foundations of Machine Learning [6h]
		\begin{itemize}
			\item Introduction
			\item Typology of Machine Learning Problems
			\item Core Issues in Artificial Learning
			\item Linear Models
			\item Validation of Machine Learning Algorithms
		\end{itemize}
		\item Machine Learning Algorithms [13h30]
		\begin{itemize}
			\item Naïve Bayes Algorithm [1h30]
			\item Principal Component Analysis [2h]
			\item Decision Trees [2h]
			\item Perceptrons and Multilayer Perceptrons [4h]
			\item K-Means [2h]
			\item K-Nearest Neighbors [2h]
			%\item Support Vector Machines [4h]%
		\end{itemize}
	\end{itemize}
\end{frame}

\subsection[Paradigms]{Paradigms}

\begin{frame}{Paradigms and Systems of Thought}
	Daniel Kahneman (2002 Nobel Prize in Economics) theorized 2 modes of functioning of the human brain:
	\begin{itemize}
		\item \textbf{Slow thinking}: conscious, logical, effortful
		\item \textbf{Fast thinking}: unconscious, intuitive, effortless
	\end{itemize}
	\vspace{0.5cm}
	\begin{columns}
		\begin{column}{0.5\textwidth}
			\begin{figure}
				\centering
				\includegraphics[height=0.5\linewidth]{Images/velo_centrifuge}
				\label{fig:velo-centrifuge}
			\end{figure}
			%\vspace{-0.5cm}
			\centering
			\small Logical approach
		\end{column}
		\begin{column}{0.5\textwidth}
			\begin{figure}
				\centering
				\includegraphics[height=0.5\linewidth]{Images/velo-t-es-capable}
				%\caption{approche intuitive}
				\label{fig:velo-t-es-capable}
			\end{figure}
			%\vspace{-0.5cm}
			\centering
			\small Intuitive approach
		\end{column}
	\end{columns}
	\vfill 
	\tiny \textcolor{gray}{Source: Thinking, fast and slow - D. Kahneman - 2012 - Penguin}
\end{frame}

\begin{frame}{Paradigms and Systems of Thought}
	A link can be established with programming paradigms:
	\begin{itemize}
		\item Slow thinking $\rightarrow$ \textbf{Deductive} programming: based on logical and explicit rules, derived from human knowledge. The program applies the rules to the provided data.
		\item Fast thinking $\rightarrow$ \textbf{Inductive} programming: based on data and experience; the rules are implicit and discovered by the algorithm. The program applies the rules it has discovered to new data; it generalizes.
	\end{itemize}
\end{frame}

%\begin{frame}{Paradigms and Systems of Thought}
%	Nevertheless, another paradigm exists:
%	\begin{itemize}
	%		\item Mimicry $\rightarrow$ \textbf{Transductive} programming: there are no rules, implicit or explicit, independent of the training data. The algorithm does not generalize, it copies.
	%	\end{itemize}
%	\begin{figure}
	%		\centering
	%		\includegraphics[height=0.3\linewidth]{Images/mimétisme}
	%		\label{fig:mimétisme}
	%	\end{figure}
%	%\vspace{-0.5cm}
%	\centering
%	\small Mimetic approach
%\end{frame}

\subsection[Definitions]{Definitions}

\begin{frame}{Definitions of Machine Learning}
	\begin{block}{Definition 1}
		Field of study that gives computers the ability to learn without being explicitly programmed. \\
		\hfill
		\textit{--- Arthur Samuel (1959)}
	\end{block}
	\vfill
	\tiny \textcolor{gray}{Source: Some Studies in Machine Learning Using the Game of Checkers - A.L. Samuel - 1959 - IBM Journal}
\end{frame}

\begin{frame}{Definitions of Machine Learning}
	\begin{block}{Definition 2}
		A program learns from \textbf{experience $E$} for a \textbf{task $T$} with a \textbf{performance $P$}, if $P$ improves with $E$. \\
		\hfill 
		\textit{--- Tom Mitchell (1997)}
	\end{block}
	%\pause
	\begin{itemize}
		\item \textbf{Experience ($E$)} = Data.
		\item \textbf{Task ($T$)} = Predict, classify, cluster, act, generate.
		\item \textbf{Performance ($P$)} = The mathematical loss function.
	\end{itemize}
	\vfill 
	\tiny \textcolor{gray}{Source: Machine Learning - T.M. Mitchell - 1997 - McGraw-Hill}
\end{frame}

\begin{frame}{Everything is a number!}
	\begin{figure}
		\centering
		\includegraphics[width=1.0\linewidth]{Images/tout_est_nombre}
		\label{fig:tout-est-nombre}
	\end{figure}
\end{frame}